\section{Introduction}

Self-supervised learning has been the most successful approach to semi-supervised learning, leveraging unlabeled data for various downstream tasks \citep{zhang2022bigssl}.
The impact of self-supervised learning in speech processing has now been extended to the pursuit of more accessible yet compact, discrete representations to interact with language models \citep{lakhotia2021generative,borsos2023audiolm,wang2023neural}.
Behind these successes, other than the ease of scaling with Transformers, the training objectives are the major factor that makes the advancement possible.

The training objective of HuBERT \citep{hsu2021hubert} is undoubtedly the most successful self-supervised objective.
The HuBERT objective combines quantization and masked prediction as its main components.
However, those choices are supported by empirical evidence rather than an underlying principle, and it becomes difficult to further improve the HuBERT objective and to design novel objectives.
This is evidenced by the fact that subsequent work has focused on data augmentation (e.g., WavLM \citep{chen2022wavlm}), simplification of training (e.g., BEST-RQ \citep{chiu2022self} and MelHuBERT \citep{lin2022melhubert}), 
pairing with other objectives (e.g., w2v-BERT \citep{chung2021w2v}, DinoSR \citep{liu2024dinosr}, 
and MS-HuBERT \citep{yadav2024ms}), 
training with different resolution \citep{shi2023multi}; none has improved the HuBERT objective itself.

In this paper, we will show that predictive coding is the underlying principle of HuBERT.
It is not difficult to see the lineage of predictive coding in HuBERT.
HuBERT is inspired by wav2vec 2.0 \citep{baevski2020wav2vec} and DeepCluster \citep{caron2018deep}, 
and wav2vec 2.0 is in turn influenced by BERT \citep{devlin2018bert} and contrastive predictive coding (CPC) \citep{oord2018representation}.
However, the HuBERT objective looks distinctively different from, say, the formulation in \citet{atal1971speech}, \citet{srinivasan1982predictive}, and \citet{rao1999predictive}.
What is missing is a general framework for predictive coding (to subsume masked prediction) and a discrete intermediate representation (to subsume quantization).

In this paper, the general framework we develop is a variational view of predictive coding and we present HuBERT as a special case.
We will derive a training objective from first principles, and will show how it relates, not only to the HuBERT objective but also to various other objectives, such as VQ-APC \citep{chung2020vector}, VQ-CPC \citep{niekerk2020vector}, wav2vec 2.0, and BEST-RQ.
Moreover, our derivation can spawn new objectives, and we will give an example that brings immediate improvement over HuBERT.
Once we have an objective, the optimization of it, i.e., the algorithm of learning representations, naturally decouples, and provides opportunities for improvement.
We will give an example of how optimizing the HuBERT objective with a different algorithm brings immediate improvement over HuBERT.
We will also validate how much the improvement on the self-supervised objective transfers to downstream performance.

%
%
%
%
%
%
%
%
%
%
%
%
%

%
%
%
